\subsection{全体概要：この章で学ぶこと}
この章は、ソフトウェアを作り、運用し、保守し、最後に廃止するまでの「作業の組み立て方」を扱う。ここでいうプロセスは、プログラムが実行中に進む処理の流れではなく、人やチームが成果物を作るために行う作業活動である。

この章の中心は三つである。第一に、プロセスとは何かという基礎概念である。第二に、複数のプロセスを時間軸や反復構造の中に配置するライフサイクルである。第三に、実行中のプロセスを測定し、評価し、改善する方法である。

初学者は、ウォーターフォール、アジャイル、DevOps といった名前を暗記する前に、「入力を活動で変換して出力を作る」「出力は別の活動の入力になる」「すべてのプロジェクトに万能なプロセスはない」「測定なしにプロセスの良し悪しは判断しにくい」という四点を押さえるとよい。

\begin{pointbox}{到達目標}この章を読み終えると、プロセス、活動、タスク、ライフサイクル、ライフサイクルモデル、プロセス評価、プロセス改善の違いを説明できる。また、予測型、反復型、進化型、増分型、継続的開発、アジャイル、DevOps を、単なる流行語ではなく、要求の変化、納品頻度、測定、組織連携という観点から比較できるようになる。\end{pointbox}
\par\smallskip
\begin{center}
\centering
\IfFileExists{assets/generated_figures/ch10/fig-ch10-01.png}{\includegraphics[width=0.96\linewidth]{assets/generated_figures/ch10/fig-ch10-01.png}}{\fbox{\parbox[c][48mm][c]{0.9\linewidth}{\centering 画像生成待ち\\10章の全体地図}}}
\par\smallskip
\refstepcounter{figure}\label{fig:ch10-01}図 \thefigure: 10章の全体地図
\end{center}
図 \thefigure は、10章の全体地図を視覚的に整理したものである。
\par\smallskip
